\section{The Rights of Man: Life}

Ever since the French Revolution, there has been much discussed about the so-called rights of man. Over time, this discussion has confounded natural rights with civil rights, i.e., the rights accorded by law to a citizen, and goods, such as food, medicine, education, which are the products of human activity. In order to make sense of rights, they must be related to the virtue of Justice, the highest virtue as it is the manifestation of the cosmic order. The principle of justice is ``just deserts'', i.e., ``to each what he deserves''.

There are three fundamental rights---the rights to life, liberty, and property---that derive from the very nature of man and hence are unalienable and do not depend, as a matter of justice, on any grant or law of any government. Since it is the task of man to know, love, and serve, these rights are necessary for the fulfillment of those tasks.

\paragraph{Life}


Life, along with space, time, form, and matter, is one of the fundamental conditions for manifestation\footnote{\url{https://gornahoor.net/?p=4927}}. In particular, this means that life does not derive from matter, a lifeless universe is impossible, and that there is an interior aspect to things along with the sensual. Hence, the Thomist principle that ``all knowledge beings with the senses'' must be balanced with the Augustinian maxim that ``truth resides in the interiority of man''.

That is because life is the same as soul, both derived from the word ``anima'' for soul and a living being is ``animated''. This teaching has been expressed as the ``anima mundi'', with variations of it in most traditions, including the Hermetic. \textbf{Valentin Tomberg}, following
\textbf{Vladimir Solovyov}, understands this in a deeper sense as

Sophia, the Eternal Feminine, or Divine Wisdom. Thus, the Intellect is feminine and reflects the Love which acts. This is the principle of knowing.

For a physical thing, our knowing begins with its sense impressions, i.e., its material reflection in the soul. However, it is not yet knowledge until its essence is known. At that point, we hold in the mind its idea, the same idea that also exists in the divine mind. Hence, we are the thing, at least in its essence although not in existence, since we participate in the form. Thomism is the philosophical system that derives from this fundamental insight by its logical proofs.

Nevertheless, for the Hermetic method, this is only the starting point and it must be supplemented with the awareness of inner states, not just external objects. Correctly understood, and we have seen this in our selected passages from \textbf{Julius Evola\footnote{\url{https://gornahoor.net/?p=6261}}}, hylomorphism has a more general applicability than to just physical things. ``Matter'' in the sense that it is ``informed'' by an essence transcends physicality. Hence, we are justified in using our impressions of extended matter as a takeoff point.

In particular, the Delphic maxim to ``know thyself'' is the imperative of the right to life. For our deepest understanding does not arise from the study of physical objects, but rather from the self-knowledge of our own soul life. In addition to the five outer senses, there are the five inner senses, or wits\footnote{\url{https://gornahoor.net/?p=3559}}. These latter inner senses are likewise the starting points for knowledge.

Beyond them, and deeper, there are the various emotions and then thinking itself. These are also sensations that can be observed and known. Beings, such as angels, are examples of formless manifestation, and thus are not known in the same way as material things. Although they sometimes can be experienced through the inner wit of fantasy, normally they are experienced as thought forms. For example, a creative idea or an uplifting thought may be such experiences. Once again, to know is to be, so to know that angel is to participate in the same idea or form, which can also be understood as a higher state of being.

Clearly, to reach these higher states, there must be the possibility to do so, viz., man has the potency to transcend his ``normal'' state. Although in a philosophical system, to ``know God'' means to know certain propositions or properties about him. However, in the Hermetic and mystical sense, ``to know'' means a direct intuition, unmediated by logical discourse. Hence, as \textbf{Louis-Claude de Saint-Martin\footnote{\url{https://www.meditationsonthetarot.com/the-tableau-of-god}} writes}:

\begin{quotex}
The only means that we would therefore have of proving the true God, the God ruling over free beings, finally the loving God and source of a joy that can be communicated to other beings, would be without doubt to demonstrate in his creature the existence of some base or some essence similar to his, and even to obtain and to feel this joy of which He is the principle.
\end{quotex}


The knowledge of the living God is obviously quite different from the knowledge of the ``logical God'' of the philosophers, as sound as that latter knowledge may be. This is the solitary path of man, since it depends on one's interior recognition of the presence of God. But, as Saint-Martin points out, ``to recognize a thing, is not always to demonstrate it.''

To deny this is to deny God himself. Since the Intellect is the reflection of the will of God, Saint-Martin's claim makes sense, namely that in the soul ``we can see reflected, as in a mirror, all the pure and sacred rays, of which the tableau of God must be composed.''

Now, one of the inner wits is called ``estimation'', which is still used when people say, ``in my estimation ... '' This is more than mere opinion, not exactly an instinct, but rather a form of natural knowledge. The soul, or life, is prior to thought, hence there is a Sophia, or Wisdom, which is certain knowledge, since it is the reflection of God's tableau, even before it can be articulated in thought. This estimation is our natural feelings of our animal faith in the existence of things, our love for family and nation, the rule of the fathers, hierarchy, moral principles, the joy of fertility, the right to free association, and so on.

These are based on a trust in natural goodness and the accumulated wisdom of human life that lies behind customs, mores, traditions, and prejudices. The truly rational man may try to make these things clear and support them with sound reasons. However, the forces of revolution always attack this accumulated wisdom as nothing more than arbitrary, and hence unjust, biases. When men lose touch with their inner life force, they are

\begin{quotex}
tossed to and fro, and carried about with every wind of doctrine by the wickedness of men, by cunning craftiness, by which they lie in wait to deceive. 
\flright{\textsc{Ephesians 4:14}}
\end{quotex}


Thus, revolution always begins with an attack on the right to life, first corporeally and then in the soul life of a people.

To protect the right to life, the political power needs to enforce several derivative civil rights. These would include the rights to self defense, to stand one's ground, to free association, to educate one's children (this is not the same as the pseudo-right to government run education), to police protection, and so on.


\flrightit{Posted on 2013-09-04 by Cologero}

